\chapter{System Design and Architecture}

    \section{Use Case Diagram}
    The Use Case Diagram shows the interactions between the student (primary user) and the system functionalities. The primary actor (Student) can perform the following use cases:
    \begin{itemize}
        \item Register and Login
        \item Manage Tasks (Create, Update, Delete, Complete)
        \item Schedule Study Sessions
        \item Run Focus Timer (Pomodoro)
        \item Take Quizzes
        \item Write Notes and Journal Entries
        \item View Progress and Analytics
        \item Export Reports
    \end{itemize}
    \begin{figure}[htbp]
        \centering
        \includegraphics[width=0.92\textwidth]{img/finalusecase.png}
        \caption{Use Case Diagram of StudyBug}
    \end{figure}

    \clearpage
    \section{Context Diagram}
    The Context Diagram shows the top level view of the StudyBug system and its interaction with external entities.
    \begin{figure}[H]
        \centering
        \includegraphics[width=1.0\textwidth]{img/contextdiagram.png}
        \caption{Context Diagram of StudyBug}
    \end{figure}

    \clearpage
    \section{Sequence Diagram}
    The Sequence Diagram illustrates the flow of a typical focus session workflow.
    \begin{figure}[H]
        \centering
        \includegraphics[width=1.0\textwidth]{img/sequence-diagram}
        \caption{Sequence Diagram - Focus Session Workflow}
    \end{figure}

    \clearpage
    \section{Data Flow Diagram}
    The Data Flow Diagram shows the flow of data between the subsystems of StudyBug.
    \begin{figure}[H]
        \centering
        \includegraphics[width=0.92\textwidth]{img/dataflow.png}
        \caption{Data Flow Diagram of StudyBug}
    \end{figure}

    \section{System Block Diagram}
    The system block diagram illustrates the architecture of StudyBug showing how different modules interact.
    \begin{figure}[H]
        \centering
        \includegraphics[width=0.82\textwidth]{img/System-Block-Diagram.png}
        \caption{System Block Diagram of StudyBug}
        \label{fig:System Block Diagram}
    \end{figure}

    \subsection{Frontend Interface Module}
    The frontend interface module handles user interaction and visual presentation. Developed using React 19, it provides interfaces for:
    \begin{itemize}
        \item Dashboard with learning progress
        \item Virtual Study Room with interactive hotspots
        \item Task management and calendar
        \item Notes and journal pages
        \item Quiz interface
        \item Focus timeline analytics
    \end{itemize}

    \subsection{Authentication Module}
    The authentication module manages user login and access control using Supabase Auth. It verifies credentials and issues JWT tokens for secure communication between frontend and backend.

    \subsection{Backend Server Module}
    The Express.js backend server acts as the central control unit, handling:
    \begin{itemize}
        \item RESTful API endpoints (35+ endpoints)
        \item Request validation and error handling
        \item Business logic processing
        \item Database communication
    \end{itemize}

    \subsection{Database Module}
    PostgreSQL database stores all persistent data across 12 tables:
    \begin{itemize}
        \item users, tasks, notes, journal\_entries
        \item calendar\_events, files, focus\_sessions
        \item learning\_sections, learning\_levels
        \item progress, trash, schedules
    \end{itemize}

    \clearpage
    \section{Entity Relationship Analysis}
    The Entity-Relationship analysis defines the data architecture of StudyBug.

    \subsection{Entity Classification}
    \begin{itemize}
        \item \textbf{Strong Entities}:
        \begin{itemize}
            \item \textbf{USER}: Primary actor; identified by \texttt{id}
            \item \textbf{TASK}: Work unit created by user; identified by \texttt{task\_id}
            \item \textbf{SCHEDULE}: Planned time allocations; identified by \texttt{schedule\_id}
            \item \textbf{NOTE}: Quick notes; identified by \texttt{note\_id}
            \item \textbf{JOURNAL\_ENTRY}: Daily entries; identified by \texttt{entry\_id}
            \item \textbf{FILE}: Uploaded documents; identified by \texttt{file\_id}
        \end{itemize}
        \item \textbf{Weak Entities}:
        \begin{itemize}
            \item \textbf{PROGRESS}: Dependent on TASK, tracks task state
            \item \textbf{FOCUS-SESSION}: Dependent on TASK, logs execution instances
            \item \textbf{LEARNING\_LEVEL}: Dependent on LEARNING\_SECTION
        \end{itemize}
    \end{itemize}

    \subsection{Relationships and Cardinality}
    \begin{table}[htbp]
    \centering
    \small
    \begin{tabular}{|p{3.8cm}|p{3.8cm}|c|p{4.5cm}|}
    \hline
    \textbf{Entity A} & \textbf{Entity B} & \textbf{Cardinality} & \textbf{Description} \\ \hline
    USER & TASK & $1:N$ & One user can manage multiple tasks \\ \hline
    USER & NOTE & $1:N$ & One user can create multiple notes \\ \hline
    USER & JOURNAL\_ENTRY & $1:N$ & One user can have multiple journal entries \\ \hline
    USER & FILE & $1:N$ & One user can upload multiple files \\ \hline
    TASK & FOCUS-SESSION & $1:N$ & A task can have many focus sessions \\ \hline
    TASK & CALENDAR\_EVENT & $1:N$ & A task can be scheduled multiple times \\ \hline
    TASK & PROGRESS & $1:N$ & A task generates progress records over time \\ \hline
    USER & LEARNING\_SECTION & $1:N$ & One user has multiple learning sections \\ \hline
    LEARNING\_SECTION & LEARNING\_LEVEL & $1:N$ & Each section has multiple levels \\ \hline
    \end{tabular}
    \caption{Entity Relationship and Cardinality Table}
    \end{table}

    \begin{figure}[H]
        \centering
        \includegraphics[width=0.9\textwidth]{img/erd2.png}
        \caption{ER Diagram of StudyBug}
    \end{figure}

    \begin{figure}[H]
        \centering
        \includegraphics[width=0.9\textwidth]{img/erd1.png}
        \caption{Detailed Entity Relationship Diagram}
    \end{figure}

    \clearpage
    \section{API Architecture}
    The backend exposes RESTful API endpoints organized by resource:

    \begin{table}[H]
    \centering
    \begin{tabular}{|p{4cm}|p{5cm}|p{4.5cm}|}
    \hline
    \textbf{Route} & \textbf{Methods} & \textbf{Purpose} \\
    \hline
    /api/auth & POST & Authentication (login, register, logout) \\
    \hline
    /api/tasks & GET, POST, PUT, DELETE & Task management \\
    \hline
    /api/notes & GET, POST, PUT, DELETE & Notes management \\
    \hline
    /api/journal & GET, POST, PUT, DELETE & Journal entries \\
    \hline
    /api/calendar & GET, POST, PUT, DELETE & Calendar events \\
    \hline
    /api/files & GET, POST, DELETE & File management \\
    \hline
    /api/focus-sessions & GET, POST, PUT & Focus session tracking \\
    \hline
    /api/progress & GET, PUT & Learning progress \\
    \hline
    /api/quiz & GET, POST & Quiz generation and results \\
    \hline
    /api/trash & GET, POST, DELETE & Soft delete management \\
    \hline
    \end{tabular}
    \caption{API Route Architecture}
    \end{table}
